\section{The de Rham complex of master function} \label{sec DR}

\subsection{Twisted de Rham complex} \label{def dr}
Consider ${\mathbb C}$ with coordinate $t$. Define the {\it master function} by the formula
\begin{gather*}%\label{Master}
\Phi(t) = \prod_{i=1}^n (t-z_i)^{-m_i/\kappa} ,
\end{gather*}
where $z_1, \dots, z_n, m_1,\dots,m_n, \kappa$ $\in{\mathbb C}$ are parameters. Fix these parameters and assume that $z_1,\dots,z_n$ are distinct. Set
\begin{gather*}
z_{n+1}=\infty, \qquad m_{n+1}=m_1+\dots+m_n-2.
\end{gather*}
Denote $U={\mathbb C}-\{z_1,\dots,z_n\}$.

Consider the {\em twisted de Rham complex} associated with $\Phi$,
\begin{gather}\label{dr com}
0\longrightarrow \Omega^0(U)\overset{\partial}{\longrightarrow}\Omega^1(U) \longrightarrow 0.
\end{gather}
Here $\Omega^p(U)$ is the space of rational differential $p$-forms on ${\mathbb C}$ regular on $U$. The differential $\partial$ is given by the formula
\begin{gather}\label{nabla}
\partial={\rm d}+\alpha\wedge\cdot ,
\end{gather}
where ${\rm d}$ is the standard de Rham differential and the second summand is the left exterior multiplication by the form
\begin{gather*}%\label{alpha}
\alpha=-\frac{1}{\kappa}\sum_{i=1}^n m_i \frac{{\rm d}t}{t-z_i} =\frac{{\rm d}\Phi}\Phi.
\end{gather*}
Formula (\ref{nabla}) is motivated by the computation
\begin{gather*}
{\rm d}(\Phi\omega)=\Phi {\rm d}\omega+{\rm d}\Phi\wedge\omega=\Phi({\rm d}\omega+\alpha\wedge\omega).
\end{gather*}
The complex $\Omega^{\bullet}(U)$ is the complex of global {\em algebraic} sections of
the de Rham complex of $\big({\mathcal{O}}^{\rm an}_U,\partial\big)$, where $\partial={\rm d}+\alpha\wedge\cdot$ is considered as the integrable connection on the sheaf ${\mathcal{O}}^{\rm an}_U$ of holomorphic functions on $U$.

If the monodromy of $\Phi$ is non-trivial, that is, if at least one of the numbers $m_1/\kappa, \dots, m_n/\kappa$ is not an integer, then
\begin{gather*}
H^0(\Omega^{\bullet}(U))=0,\qquad \dim H^1(\Omega^{\bullet}(U))=n-1,
\end{gather*}
see for example \cite{STV}.

\subsection[Basis of $\Omega^{\bullet}(U)$]{Basis of $\boldsymbol{\Omega^{\bullet}(U)}$} \label{Basis of Omega}

The functions
\begin{gather*}
\frac1{(t-z_i)^{a}}\ \ \text{for}\ a\in{\mathbb Z}_{>0} \qquad\text{and}\qquad t^a\ \ \text{for}\ a\in{\mathbb Z}_{\geq 0}
\end{gather*}
form a basis of $\Omega^0(U)$. The differential forms
\begin{gather*}
\frac{{\rm d}t}{(t-z_i)^a}\ \ \text{for}\ a\in{\mathbb Z}_{>0} \qquad \text{and}\qquad t^a{\rm d}t\ \ \text{for}\ a\in{\mathbb Z}_{\geq 0}
\end{gather*}
form a basis of $\Omega^1(U)$. The differential $\partial$ is given by the formulas
\begin{gather}
\kappa \partial\left(\frac1{(t-z_i)^{a}}\right) = -(m_i+a\kappa)\frac{{\rm d} t}{(t-z_i)^{a+1}}
+ \sum_{k=1}^a\sum_{j\neq i} \frac{m_j}{(z_j-z_i)^k} \frac{{\rm d}t}{(t-z_i)^{a+1-k}} \nonumber\\
\hphantom{\kappa \partial\left(\frac1{(t-z_i)^{a}}\right) =}{} -\sum_{j\neq i}\frac{ m_j}{(z_j-z_i)^a} \frac{{\rm d}t}{t-z_j}, \label{bff} \\
\kappa \partial\big(t^a\big) = \left(a\kappa -\sum_{j=1}^n m_j\right)t^{a-1}dt -\sum_{k=1}^{a-1} \sum_{j=1}^n m_jz_j^k t^{a-1-k}{\rm d}t-\sum_{j=1}^n m_jz_j^a \frac{{\rm d}t}{t-z_j}.\label{bfi}
\end{gather}


\subsection{Resonances}\label{reson}
The equations
\begin{enumerate}\itemsep=0pt
\item[(i)] $m_i+(a-1)\kappa=0$ for some $a\in {\mathbb Z}_{>0}$, $i\in\{1,\dots,n\}$,
\item[(ii)] $m_{n+1}+2-a\kappa=0$ for some $a\in {\mathbb Z}_{>0}$,
\item[(iii)] $\kappa=0$,
\end{enumerate}
are called the {\it resonance relations} for the parameters $m_1,\dots,m_{n+1}$, $\kappa$ of the de Rham complex.

If $\kappa=0$, then the twisted de Rham complex is not defined. If the resonance relation \mbox{$m_i+a\kappa=0$} is satisfied for some $a$, then the first term in the right-hand side of (\ref{bff}) equals zero. Similarly, if the resonance relation $m_{n+1}+2-a\kappa=0$ is satisfied for some $a$, then the first term in the right-hand side of (\ref{bfi}) equals zero.

\subsection{Logarithmic subcomplex} \label{log forms}
Let $\Omega_{\log}^0(U) \subset \Omega^0(U)$ be the subspace generated over ${\mathbb C}$ by function $1$. Let $\Omega_{\log}^1(U) \subset \Omega^1(U)$ be the subspace generated over ${\mathbb C}$ by the differential forms
\begin{gather*}
\omega_j=\frac{{\rm d}t}{t-z_j}, \qquad j=1,\dots,n.
\end{gather*}
These subspaces form the {\it logarithmic} subcomplex $(\Omega^\bullet_{\log}(U), \partial)$ of the de Rham complex\linebreak $(\Omega^\bullet(U), \partial)$. We have
\begin{gather*}
\partial\colon \ 1 \mapsto \alpha.
\end{gather*}

For generic $m_1,\dots, m_n$, $\kappa$, the embedding $(\Omega^\bullet_{\log}(U), \partial)\hookrightarrow (\Omega^\bullet(U), \partial)$ is a quasi-isomorphism, the logarithmic forms $\omega_1,\dots,\omega_n$ generate the space $H^1(\Omega^\bullet(U))$, and the cohomological relation $\sum\limits_{i=1}^n m_i \omega_i\sim 0$ is the only one, see for example~\cite{STV}.

Each resonance relation implies a new cohomological relation between the forms $\omega_1,\dots,\omega_n$, see \cite[Corollary 6.4]{SV2}. For example, if $m_{n+1}+2-\kappa=0$, then $\sum\limits_{j=1}^n z_j m_j\omega_j \sim 0$, and if $m_{n+1}+2-2\kappa=0$, then
\begin{gather*}
\sum_{j=1}^n z_j^2m_j\omega_j-\frac{1}{\kappa}\left(\sum_{j=1}^n\ z_jm_j\right)\left(\sum_{i=1}^n\ z_im_i\omega_i\right)\sim 0.
\end{gather*}

\section[$\widehat{{\mathfrak{sl}_2}}$-modules]{$\boldsymbol{\widehat{{\mathfrak{sl}_2}}}$-modules} \label{sec hsl}
\subsection[Lie algebra $\widehat{{\mathfrak{sl}_2}}$]{Lie algebra $\boldsymbol{\widehat{{\mathfrak{sl}_2}}}$} \label{lie alg hsl}

Let $\mathfrak{sl}_2$ be the Lie algebra of complex $(2\times 2)$-matrices with zero trace. Let $e$, $f$, $h$ be standard generators subject to the relations
\begin{gather*}
[e,f]=h,\qquad [h,e]=2e,\qquad [h,f]=-2f.
\end{gather*}
Let $\widehat{{\mathfrak{sl}_2}}$ be the affine Lie algebra $\widehat{{\mathfrak{sl}_2}}=\mathfrak{sl}_2\big[T,T^{-1}\big]\oplus{\mathbb C} c$ with the bracket
\begin{gather*}
\big[aT^i,bT^j\big] = [a,b]T^{i+j} + i \langle a,b\rangle \delta_{i+j,0} c ,
\end{gather*}
where $c$ is central element, $\langle a,b\rangle=\operatorname{tr} (ab)$. Set
\begin{alignat*}{4}
&e_1 = e, \qquad && f_1 = f, \qquad && h_1=h,& \\
& e_2 = fT, \qquad && f_2=eT^{-1}, \qquad && h_2=c-h.&
\end{alignat*}
These are the standard Chevalley generators defining $\widehat{{\mathfrak{sl}_2}}$ as the Kac--Moody algebra correspon\-ding to the Cartan matrix
$\left(\begin{smallmatrix} \hphantom{-}2&-2 \\
-2&\hphantom{-}2 \end{smallmatrix}\right).$

\subsection[Automorphism $\pi$]{Automorphism $\boldsymbol{\pi}$} \label{aut}
The Lie algebra $\widehat{{\mathfrak{sl}_2}}$ has an automorphism $\pi$,
\begin{gather*}
\pi \colon \ c \mapsto c, \qquad eT^i \mapsto fT^i, \qquad fT^i\mapsto eT^i, \qquad hT^i\mapsto -hT^i.
\end{gather*}

\subsection{Verma modules} \label{Ver mod}
We fix $k\in{\mathbb C}$ and assume that the central element~$c$ acts on all our representations by multiplication by~$k$.

For $m\in{\mathbb C}$, let $V(m,k-m)$ be the $\widehat{{\mathfrak{sl}_2}}$ {\em Verma module} with generating vector $v$. The Verma module is generated by $v$ subject to the relations
\begin{gather*}
e_1v=0,\qquad e_2v=0,\qquad h_1v=mv,\qquad h_2v=(k-m)v.
\end{gather*}

Let $\hat{\mathfrak{n}}_-\subset\widehat{{\mathfrak{sl}_2}}$ be the Lie subalgebra generated by $f_1$, $f_2$ and $U\hat{\mathfrak{n}}_-$ its enveloping algebra. The map $U\hat{\mathfrak{n}}_-\to V(m,k-m)$, $F\mapsto Fv$, is an isomorphism of $U\hat{\mathfrak{n}}_-$-modules.

The space $V(m,k-m)$ has a ${\mathbb Z}^2_{\geq 0}$-grading: a vector $f_{i_1}\cdots f_{i_p}v$\ with $i_j\in\{1,2\}$ has deg\-ree~$(p_1,p_2)$, if $p_i$ is the number of $i$'s in the sequence $i_1,\dots,i_p$. For $\gamma\in{\mathbb Z}^2_{\geq 0}$, denote by $V(m,k-m)_{\gamma}\subset V(m,k-m)$ the corresponding $\gamma$-homogeneous component.

A homogeneous nonzero vector $\omega$ in $V(m,k-m)$, non-proportional to $v$, is called a {\em singular vector} if $e_1\omega=e_2\omega=0$. The Verma module $V(m,k-m)$ is reducible, if and only if it contains a singular vector.

\subsection{Reducibility conditions}\label{reduc}
See Kac--Kazhdan \cite{KK}. Set
\begin{gather*}
\kappa=k+2.
\end{gather*}
The Verma module $V(m,k-m)$ is reducible if and only if at least one of the following relations holds:
\begin{enumerate}\itemsep=0pt
\item[(a)] $m-l+1+(a-1)\kappa=0$,
\item[(b)] $m+l+1-a\kappa=0$,
\item[(c)] $\kappa=0$,
\end{enumerate}
where $l,a\in{\mathbb Z}_{>0}$. If $(m,\kappa)$ satisfies exactly one of the conditions~(a),~(b), then $V(m,k-m)$ contains a unique proper submodule, and this submodule is generated by a~singular vector of degree $(la,l(a-1))$ for condition~(a) and of degree $(l(a-1),la)$ for condition~(b).

These singular vectors are highly nontrivial and are given by the following theorem.

\begin{thm}[{Malikov--Feigin--Fuchs, \cite{MFF}}]\label{thm MFF} For $a, l\in{\mathbb Z}_{>0}$ and $\kappa\in {\mathbb C}$, the monomials
\begin{gather*}
F_{12}(l,a,\kappa) = f_1^{l+(a-1)\kappa}f_2^{l+(a-2)\kappa} f_1^{l+(a-3)\kappa}\cdots f_2^{l-(a-2)\kappa} f_1^{l-(a-1)\kappa}, \\ %\label{a} \\
F_{21}(l,a,\kappa) = f_2^{l+(a-1)\kappa}f_1^{l+(a-2)\kappa} f_2^{l+(a-3)\kappa}\cdots f_1^{l-(a-2)\kappa} f_2^{l-(a-1)\kappa} \nonumber
\end{gather*}
are well-defined as elements of $U\hat{\mathfrak{n}}_-$. If $m=l-1-(a-1)\kappa$, then $F_{12}(l,a,\kappa)v\in V(m,k-m)$ is a singular vector of degree $(la,l(a-1))$ and if $m=-l-1+a\kappa$, then $F_{21}(l,a,\kappa)v\in V(m,k-m)$ is a singular vector of degree $(l(a-1),la)$.
\end{thm}

An explanation of the meaning of complex powers in these formulas see in~\cite{MFF}.

For example for $m=-2+\kappa$, we have
\begin{gather*}
F_{21}(1,1,\kappa)v=f_2v=\frac{e}{T}v,
\end{gather*}
and for $m=-2+2\kappa$, we have
\begin{gather*}
F_{21}(1,2,\kappa)v=f_2^{1+\kappa} f_1f_2^{1-\kappa}v=f\left(\frac{e}{T}\right)^2 v+(1+\kappa)\frac{h}{T} \frac{e}{T}v-(1+\kappa)\kappa\frac{e}{T^2}v.
\end{gather*}

\subsection{Shapovalov form}\label{contra}
The {\em Shapovalov form} on an $\widehat{{\mathfrak{sl}_2}}$ Verma module $V$ with generating vector $v$ is the unique symmetric bilinear form $S(\cdot,\cdot)$ on $V$ such that
\begin{gather*}
S(v,v)=1,\qquad S(f_ix,y)=S(x,e_iy) \qquad \text{for} \ i=1,2;\ x,y\in V.
\end{gather*}

For $\gamma\in{\mathbb Z}^2_{\geq 0}$, let $V^*_{\gamma}$ be the vector space dual to $V_{\gamma}$. Define $V^*=\oplus_\gamma V_{\gamma}^*$. The space $V^*$ is an $\widehat{{\mathfrak{sl}_2}}$-module with the $\widehat{{\mathfrak{sl}_2}}$-action defined by the formulas:
\begin{gather*}
\langle f_i\phi,x\rangle=\langle \phi,e_ix\rangle, \qquad \langle e_i\phi,x\rangle=\langle\phi,f_ix\rangle ,
\end{gather*}
where $\phi\in V^*$, $x\in V$, $i=1,2 $. The $\widehat{{\mathfrak{sl}_2}}$-module $V^*$ is called the {\em contragradient} Verma module.

The Shapovalov form $S$ considered as a map $S\colon V\longrightarrow V^*$ is a morphism of $\widehat{{\mathfrak{sl}_2}}$-modules.

\subsection[Bases in $V$ and $V^*$]{Bases in $\boldsymbol{V}$ and $\boldsymbol{V^*}$}\label{BASIS}
Let $V$ be an $\widehat{{\mathfrak{sl}_2}}$ Verma module $V$. For every $\gamma=(p_1,p_2)\in{\mathbb Z}^2_{\geq 0}$ with $p_1\ne p_2$, we fix a basis in the homogeneous component $V_{\gamma}\subset V$.

For $p_1>p_2$, we fix the basis
\begin{gather*}%\label{p_1>p_2}
\left\{\frac{f}{T^{i_1}}\cdots \frac{f}{T^{i_a}} \frac{h}{T^{j_1}}\cdots\frac{h}{T^{j_b}} \frac{e}{T^{k_1}}\cdots\frac{e}{T^{k_c}}v \right\},
\end{gather*}
where
\begin{gather}
0\leq i_a\leq i_{a-1}\leq\dots\leq i_1,\quad 1\leq j_b\leq j_{b-1}\leq\dots\leq j_1,\quad 1\leq k_c\leq k_{c-1}\leq\dots\leq k_1; \nonumber\\
\sum_{s=1}^a i_s + \sum_{s=1}^b j_s + \sum_{s=1}^c k_s + a - c = p_1, \qquad \sum_{s=1}^a i_s + \sum_{s=1}^b j_s + \sum_{s=1}^c k_s = p_2.\label{indices}
\end{gather}
For $p_1 < p_2$, we fix the basis
\begin{gather*}%\label{p_1<p_2}
\left\{\frac{e}{T^{k_1}}\cdots\frac{e}{T^{k_c}} \frac{h}{T^{j_1}}\cdots\frac{h}{T^{j_b}} \frac{f}{T^{i_1}}\cdots\frac{f}{T^{i_a}}v\right\},
\end{gather*}
with the indices satisfying~(\ref{indices}). Notice that for any $x\in\mathfrak{sl}_2$ the elements $\frac x{T^i}$ and $\frac x{T^j}$ commute.

These collections of vectors are bases by the Poincar\'{e}--Birkhoff--Witt theorem.

\looseness=-1 For any $\gamma$, we fix a basis in the $\gamma$-homogeneous component $V_\gamma^*\subset V^*$ as the basis dual of the basis in $V_\gamma$ specified above.
If $\{w_i\}$ is a basis in $V_\gamma$, then we denote by $\{(w_i)^*\}$ the dual basis in~$V^*_\gamma$.

\subsection{Main formula}
\begin{thm}[{\cite[Theorem 5.12]{SV2}}]\label{thm form}
For $m, k\in{\mathbb C}$ and $a\in{\mathbb Z}_{>0}$, the following identities hold in the contragradient Verma module $V(m,k-m)^*$,
\begin{gather}
\frac{f}{T^{a-1}} (v)^* = (m+(a-1)(k+2))\left(\frac{f}{T^{a-1}}v\right)^* \nonumber \\
\hphantom{\frac{f}{T^{a-1}} (v)^* =} + \sum_{\ell=1}^{a-1}\bigg[\frac{h}{T^\ell}\left(\frac{f}{T^{a-1-\ell}}v\right)^* + 2 \frac{e}{T^\ell}\mathop{\sum_{i+j=a-1-\ell}}_{i\geq j\geq 0} \left(\frac{f}{T^i}\frac{f}{T^j}v\right)^*\bigg], \label{id A}\\
\frac{e}{T^a} (v)^*=(a(k+2)-m-2) \left(\frac{e}{T^a}v\right)^* \nonumber\\
\hphantom{\frac{e}{T^a} (v)^*=} + \sum_{\ell=0}^{a-2}\bigg[{-} \frac{h}{T^{\ell+1}} \left(\frac{e}{T^{a-\ell-1}}v\right)^* + 2 \frac{f}{T^\ell}\mathop{\sum_{i+j=a-\ell}}_{i\geq j\geq 1} \left(\frac{e}{T^i}\frac{e}{T^j}v\right)^*\bigg],\label{id B}
\end{gather}
where $v$ is the generating vector of the Verma module $V(m,k-m)$.
\end{thm}

Theorem~\ref{thm form} was announced in \cite{SV2}. The proof of Theorem~\ref{thm form} is the main result of this paper. The theorem is proved in Section~\ref{sec PROOF}.

\begin{rem} The right-hand sides of formulas~(\ref{id A}) and~(\ref{id B}) have the factors $m+(a-1)(k+2)$ and $a(k+2)-m-2$. The vanishing of these factors corresponds to the resonance conditions $m_i+(a-1)\kappa=0$ and $m_{n+1}+2-a\kappa=0$ for the de Rham complex in Section~\ref{reson}, if we recall that $\kappa=k+2$.
\end{rem}

\begin{rem}Theorem \ref{thm form} says that the action of the element $\frac{f}{T^{a-1}}$ of degree $(a,a-1)$ on the covector $(v)^*$ can be expressed in terms of the actions of the elements $\frac{h}{T^{l}}$ and $\frac{e}{T^{l}}$ of smaller degree on some other covectors. Similarly the action of the element $\frac{e}{T^{a}}$ of degree $(a-1,a)$ on the covector~$(v)^*$ can be expressed in terms of the actions of the elements $\frac{h}{T^{l}}$, $\frac{f}{T^{l}}$ of smaller degree on some other~covectors.
\end{rem}

